\section{Discussion}

\subsection{Key Finding: Propagation Without Lock-in}

Our results support a nuanced view of benchmark concentration in ML research. Citation networks strongly propagate benchmark standards---papers that cite each other share benchmarks at 24 times the rate of random pairs (Cohen's $d = 1.93$). This confirms that evaluation practices spread through social and intellectual networks.

However, we find no evidence of the ``epistemic lock-in'' that critics fear. Concentration in year $t-1$ does not predict diversity decline in year $t$ ($\beta = +1.60$, $p = 0.098$). The positive coefficient, while not significant, suggests if anything that concentration is followed by diversification rather than further concentration.

\subsection{Theoretical Implications}

These findings favor a \textbf{co-evolution model} over a \textbf{deterministic lock-in model}. In the co-evolution view, benchmark concentration and research practices evolve together without causal feedback---concentration may rise and fall based on field dynamics (new benchmark releases, paradigm shifts) rather than self-reinforcing cycles.

\subsection{Limitations}

\textbf{Small panel size.} With only 21 venue-years (18 after lag-drop), our panel regression has limited statistical power.

\textbf{Task labels as dataset proxy.} PWC task labels are coarser than specific datasets. This likely attenuates our concentration estimates.

\textbf{Observational design.} We document associations, not causal mechanisms.

\textbf{Venue scope.} Our analysis covers three general ML venues. Domain-specific venues may exhibit different dynamics.

\subsection{Broader Impact}

Claims that benchmark concentration creates inevitable ``lock-in'' appear overstated---the temporal feedback mechanism that would create such a trap is not supported by our data. The absence of self-reinforcing feedback suggests that the community can shift practices when problems become apparent.
